Moore 行列式的秩一摄动公式可由下述对角摄动公式和酉不变性推出。证明中将对所得公式求导，以确定相应的余子式矩阵。

\begin{Proposition}[Alesker，Theorem 1.1.9(ii) 与 Proposition 1.1.11]\label{ext:alesker-determinant}
设 $\mathbb H$ 是以 $1$ 为单位元、以 $(1,\mathbf i,\mathbf j,\mathbf k)$ 为基的实结合代数，其中 $\mathbf i^2=\mathbf j^2=\mathbf k^2=\mathbf i\mathbf j\mathbf k=-1$。共轭保持实系数不变，并将其余三个系数变号。四元数列向量采用右标量乘法和内积 $v^*w$，其中 $^*$ 表示共轭转置。

当 $m\ge1$ 时，每个四元数 $m\times m$ 矩阵 $A=A^*$ 都有一组四元数标准正交特征基，其特征值 $\lambda_1,\ldots,\lambda_m$ 均为实数，见 \cite{r-2} 的 Claim~1.1.7。对每个实数 $\lambda$，该列表中等于 $\lambda$ 的项数为 $\dim_{\mathbb H}\ker(A-\lambda I_m)$，故计入重数的列表在不计次序时与基的选取无关。定义 Moore 行列式为
\[
 \det A=\prod_{a=1}^m\lambda_a.
\]
于是 $\det I_m=1$。结合 Theorem~1.1.9(ii) 的酉不变性与 Example~1.1.10(a) 的实对角公式（见 \cite{r-2}），可知此乘积就是 Theorem~1.1.8 的归一化行列式。对任意四元数 $m\times m$ 矩阵 $U$，满足 $U^*U=I_m$，以及任意实对角矩阵 $T=\operatorname{diag}(t_1,\ldots,t_m)$，都有
\[
 \det(U^*AU)=\det A,\qquad
 \det(A+T)=\sum_{I\subseteq\{1,\ldots,m\}}
          \left(\prod_{i\in I}t_i\right)\det A_{I^c,I^c}.
\]
其中 $A_{I^c,I^c}$ 是在指标集 $I^c$ 上的主子矩阵，空矩阵的行列式取一。归一化及这些恒等式见 \cite{r-2}, pp.~7--8。
\end{Proposition}
